\chapter{Ergebnisse}
\label{cha:Ergebnisse}

\section{Trainingsverhalten}
\label{sec:trainingsverhalten}

Das Training wurde mit 512 parallelen Umgebungen auf einer NVIDIA A30 GPU
durchgeführt (341\,FPS Durchsatz) und nach 957~Iterationen durch die
24-Stunden-GPU-Zeitbegrenzung beendet.
Abbildung~\ref{fig:reward_training} zeigt den Belohnungsverlauf.

\begin{figure}[H]
    \centering
    \includegraphics[width=0.48\textwidth]{figures/training/Train_mean_reward.pdf}
    \hfill
    \includegraphics[width=0.48\textwidth]{figures/training/Policy_mean_std.pdf}
    \caption[Kumulative Belohnung und Standardabweichung der
    Aktionsverteilung]{Kumulative Belohnung (links) und mittlere
    Standardabweichung der Aktionsverteilung (rechts) über 957
    Trainingsiterationen. Geglättet mit EMA $\alpha = 0{,}9$,
    transparenter Hintergrund: Rohwerte.}
    \label{fig:reward_training}
\end{figure}

Der Trainingsverlauf lässt sich in zwei Phasen unterteilen. In der ersten
Phase (Iteration~0 bis~441) steigt die kumulative Belohnung kontinuierlich
an. Der Agent lernt zunächst die Zielnavigation; ab circa Iteration~200
erscheint erstmals die Erfolgsbelohnung
(Abbildung~\ref{fig:rew_crash_success}), und gleichzeitig sinkt die
Absturzrate. Die kumulative Belohnung erreicht ihren Höchstwert von 24,8
bei Iteration~390. Die Standardabweichung der Aktionsverteilung fällt
monoton von 0,99 auf 0,15, was die Konvergenz zu einer deterministischen
Strategie zeigt. Die übrigen Belohnungskomponenten bestätigen den
Lernfortschritt
(Abbildungen~\ref{fig:rew_eplen_close} und~\ref{fig:rew_prox_std} im
Anhang).

\begin{figure}[H]
    \centering
    \includegraphics[width=0.48\textwidth]{figures/training/Episode_rew_crash_raw.pdf}
    \hfill
    \includegraphics[width=0.48\textwidth]{figures/training/Episode_rew_success_raw.pdf}
    \caption[Absturzstrafe und Erfolgsbelohnung]{Unskalierte Absturzstrafe
    (links) und Erfolgsbelohnung (rechts). Geglättet mit EMA
    $\alpha = 0{,}9$.}
    \label{fig:rew_crash_success}
\end{figure}

Ab Iteration~442 bricht die kumulative Belohnung abrupt ein: Innerhalb einer
einzigen Iteration fällt sie von 29,2 auf $-6{,}1$. Die
Standardabweichung bleibt nahezu unverändert ($0{,}193 \rightarrow
0{,}194$); der Kollaps betrifft die Mittelwerte der Strategie, nicht die
Explorationsbreite.

Der beste Checkpoint liegt bei Iteration~400 (Tabelle~\ref{tab:konvergenz}).
Auf Basis des Belohnungsverlaufs wurde dieser für alle Evaluierungen
ausgewählt (\emph{Early Stopping}).

\begin{table}[H]
  \centering
  \caption[Konvergenzgeschwindigkeit und Checkpoint-Selektion]{Konvergenzgeschwindigkeit und Checkpoint-Selektion (512 Umgebungen, NVIDIA A30, 341\,FPS).}
  \label{tab:konvergenz}
  \small
  \renewcommand{\arraystretch}{1.1}
  \begin{tabular}{l r r r}
    \toprule
    \textbf{Checkpoint} & \textbf{Iteration} & \textbf{Laufzeit} & \textbf{Kum.\ Belohnung} \\
    \midrule
    Bestes Modell   & 400 & 10,4\,h & 23,0 \\
    Post-Kollaps    & 900 & 22,6\,h & $-1{,}1$ \\
    Trainingsende   & 957 & 24,0\,h & $-0{,}8$ \\
    \bottomrule
  \end{tabular}
\end{table}


\section{Phase 1: Korridornavigation}
\label{sec:phase1}

Versuchsreihe~1 (vgl.\ Abschnitt~\ref{subsec:meth_versuch1}) evaluiert die
trainierte Strategie (Checkpoint~400) über 1\,000~Episoden mit zuvor
ungesehenen Hindernisobjekten. Der Agent erreicht eine Erfolgsrate von
71,0\,\% (Abbildung~\ref{fig:phase1_outcomes}). Die fehlgeschlagenen
Episoden verteilen sich nahezu vollständig auf Hinderniskollisionen
(28,8\,\%); die verbleibenden 0,2\,\% enden durch Verlassen des
Korridors. Bodenkollisionen und Timeouts treten nicht auf.
Tabelle~\ref{tab:phase1_results} im Anhang listet alle Metriken mit
Konfidenzintervallen.

\begin{figure}[H]
    \centering
    \includegraphics[width=0.65\textwidth]{figures/eval/phase1_outcomes.pdf}
    \caption[Episodenergebnisse Phase~1]{Episodenergebnisse der
    Phase-1-Evaluation (1\,000~Episoden). Fehlerbalken zeigen
    95\,\%-Bootstrap-Konfidenzintervalle.}
    \label{fig:phase1_outcomes}
\end{figure}

Das Flugverhalten zeigt zwei wiederkehrende Muster
(Abbildung~\ref{fig:phase1_trajectories}): Bei freier Bahn vollzieht der
Agent einen annähernd vertikalen Abstieg (mediane Pfadeffizienz 1,19). In
der Nähe eines Hindernisses bremst er spät ab und weicht lateral mit
geringem Sicherheitsabstand aus (mittlerer Minimalabstand 0,52\,m).

\begin{figure}[H]
    \centering
    \includegraphics[width=0.32\textwidth]{figures/eval/trajectory_ep1.pdf}%
    \hfill
    \includegraphics[width=0.32\textwidth]{figures/eval/trajectory_ep2.pdf}%
    \hfill
    \includegraphics[width=0.32\textwidth]{figures/eval/trajectory_ep3.pdf}
    \caption[Exemplarische 3D-Trajektorien in Phase~1]{Exemplarische
    3D-Trajektorien dreier erfolgreicher Episoden. Graue Quader: AABB der
    Hindernisse. Grünes Dreieck: Startposition, roter Punkt: Landeziel.}
    \label{fig:phase1_trajectories}
\end{figure}

\subsection{Nachuntersuchung: Oberflächenbasierte Kollisionserkennung}
\label{subsec:phase1_mesh_contacts}

Das knappe Passieren bei hoher Kollisionsrate legte nahe, dass die
konservative AABB-Erkennung einen Teil der Fehlschläge verursacht. In
einer Nachuntersuchung mit 1\,000~Episoden wurde die
AABB-Kollisionserkennung durch die geometrisch exakte
Oberflächenkollisionsprüfung des Simulators ersetzt.

Mit oberflächenbasierter Kontakterkennung steigt die Erfolgsrate auf
98,4\,\% (Tabelle~\ref{tab:phase1_mesh_contacts}). Ein erheblicher
Anteil der AABB-Kollisionen sind demnach Fehlalarme: Die Drohne
durchfliegt die Begrenzungsbox, ohne die Oberfläche zu berühren.
Tabelle~\ref{tab:phase1_mesh_results} im Anhang listet alle Metriken.

\begin{table}[H]
  \centering
  \caption[Vergleich AABB- und oberflächenbasierte Kollisionserkennung]{Vergleich
  AABB-Distanz und oberflächenbasierte Kontakterkennung ($n = 1\,000$).}
  \label{tab:phase1_mesh_contacts}
  \small
  \renewcommand{\arraystretch}{1.1}
  \begin{tabular}{l r r}
    \toprule
    \textbf{Metrik} & \textbf{AABB} & \textbf{Oberfläche} \\
    \midrule
    Erfolgsrate                       & 71,0\,\%    & 98,4\,\% \\
    Hinderniskollision                & 28,8\,\%    & 1,2\,\% \\
    Verlassen des Korridors           & 0,2\,\%     & 0,2\,\% \\
    Timeout                           & 0,0\,\%     & 0,2\,\% \\
    \bottomrule
  \end{tabular}
\end{table}


\section{Phase 2: Weinberg-Transfer}
\label{sec:phase2}

Versuchsreihe~2 (vgl.\ Abschnitt~\ref{subsec:meth_versuch2}) evaluiert den
Zero-Shot-Transfer in die Weinbergumgebung über 1\,000~Episoden. Der Agent
erreicht eine Erfolgsrate von 10,4\,\%
(Abbildung~\ref{fig:phase2_outcomes}). Timeouts dominieren mit 73,7\,\%;
Terrain-Kollisionen verursachen 15,9\,\%. Der mediane Abstand zum
Landeziel am Episodenende beträgt 1,19\,m; 59,7\,\% der
Timeout-Episoden befinden sich innerhalb von 1,5\,m zum Ziel.
Tabelle~\ref{tab:phase2_results} im Anhang listet alle Metriken.

\begin{figure}[H]
    \centering
    \includegraphics[width=0.65\textwidth]{figures/eval/phase2_outcomes.pdf}
    \caption[Episodenergebnisse Phase~2]{Episodenergebnisse der
    Phase-2-Evaluation (1\,000~Episoden, Zero-Shot-Transfer).
    Fehlerbalken zeigen 95\,\%-Bootstrap-Konfidenzintervalle.}
    \label{fig:phase2_outcomes}
\end{figure}

In der Mehrzahl der Fälle erreicht der Agent die Umgebung des Landeziels,
verbleibt jedoch außerhalb des Erfolgsradius von 0,5\,m mit ineffektiven
Korrekturbewegungen.

\subsection{Sensitivitätsanalyse des Erfolgskriteriums}
\label{subsec:phase2_relaxed}

Mit einem erweiterten Erfolgsradius von 1,5\,m steigt die Erfolgsrate auf
69,7\,\% (Tabelle~\ref{tab:phase2_relaxed}). Die
Terrain-Kollisionsrate sinkt von 15,9\,\% auf 9,4\,\%, da ein Teil der
zuvor als Timeout gewerteten Episoden nun frühzeitig als Erfolg beendet
wird. Der Timeout-Anteil sinkt von 73,7\,\% auf 27,7\,\%.

\begin{table}[H]
  \centering
  \caption[Sensitivitätsanalyse des Erfolgsradius in Phase~2]{Sensitivitätsanalyse
  des Erfolgsradius~$r$ in der Weinbergumgebung ($n = 1\,000$ je Bedingung).}
  \label{tab:phase2_relaxed}
  \small
  \renewcommand{\arraystretch}{1.1}
  \begin{tabular}{l r r}
    \toprule
    \textbf{Metrik} & \textbf{$r = 0{,}5$\,m} & \textbf{$r = 1{,}5$\,m} \\
    \midrule
    Erfolgsrate           & 10,4\,\%    & 69,7\,\% \\
    Terrain-Kollision     & 15,9\,\%    & 9,4\,\% \\
    Timeout               & 73,7\,\%    & 27,7\,\% \\
    \bottomrule
  \end{tabular}
\end{table}


\section{Vergleich RL vs.\ RRT-Connect}
\label{sec:vergleich_rl_rrt}

Der RRT-Planer übertrifft den RL-Agenten in der Erfolgsrate beider
Umgebungen, wobei der Abstand im Weinberg deutlich größer ausfällt. Im
Korridor dominieren bei beiden Verfahren Hinderniskollisionen; im Weinberg
scheitert der RL-Agent überwiegend an Timeouts (73,7\,\%), während der
RRT-Planer diese auf 7,9\,\% begrenzt. Der RL-Agent landet schneller
(30,0\,s gegenüber 42,0\,s im Korridor).\footnote{Einzelne
Pfadeffizienzwerte knapp unter 1,0 sind ein Implementierungsartefakt:
Das letzte Flugsegment wird aufgrund des automatischen Zurücksetzens der
Umgebung nicht erfasst.} Dem steht ein einmaliger Trainingsaufwand von
circa 24~Stunden gegenüber. Detaillierte Metriken mit
Konfidenzintervallen finden sich in den
Anhangstabellen~\ref{tab:phase1_results} und~\ref{tab:phase2_results}. Tabelle~\ref{tab:vergleich_rl_rrt} stellt die Ergebnisse beider Verfahren
gegenüber.

\begin{table}[H]
  \centering
  \caption[Vergleich RL-Agent und RRT-Connect]{Vergleich von RL-Agent und
  RRT-Connect über beide Versuchsreihen ($n = 1\,000$ je Bedingung).}
  \label{tab:vergleich_rl_rrt}
  \footnotesize
  \renewcommand{\arraystretch}{1.15}
  \begin{tabular}{l r r r r}
    \toprule
    & \multicolumn{2}{c}{\textbf{Phase~1 (Korridor)}}
    & \multicolumn{2}{c}{\textbf{Phase~2 (Weinberg)}} \\
    \cmidrule(lr){2-3} \cmidrule(lr){4-5}
    \textbf{Metrik}
      & \textbf{RL} & \textbf{RRT}
      & \textbf{RL} & \textbf{RRT} \\
    \midrule
    Erfolgsrate
      & 71,0\,\%       & 82,0\,\%
      & 10,4\,\%       & 85,7\,\% \\
    \quad(relaxed $r = 1{,}5$\,m)
      & \multicolumn{2}{c}{n.\,a.}
      & 69,7\,\%       & 98,4\,\% \\
    \addlinespace
    Hindernis-/Terrain-Koll.
      & 28,8\,\%       & 18,0\,\%
      & 15,9\,\%       & 7,1\,\% \\
    Timeout
      & 0,0\,\%        & 0,0\,\%
      & 73,7\,\%       & 7,9\,\% \\
    \addlinespace
    Landezeit (Median)
      & 30,0\,s         & 42,0\,s
      & 36,0\,s         & 42,0\,s \\
    Pfadeffizienz (Median)
      & 1,19            & 1,09
      & 1,07            & 0,99 \\
    Min.\ Abstand (Mittel)
      & 0,52\,m         & 0,48\,m
      & 1,02\,m         & 0,46\,m \\
    \bottomrule
  \end{tabular}
\end{table}